% ============================================================
%  Muhebullah Karimzada — Teaching Dossier
%  Teaching-Stream Assistant Professor, Economics
%  Compiled with: pdflatex (run twice for cross-references)
%  Engine: pdfLaTeX (TeX Live 2025)
%  Last updated: September 2026
% ============================================================

\documentclass[11pt,letterpaper]{article}

\usepackage[T1]{fontenc}
\usepackage[utf8]{inputenc}
\usepackage[sc]{mathpazo}
\usepackage{microtype}
\usepackage[top=1.0in, bottom=1.0in, left=1.1in, right=1.1in]{geometry}
\usepackage{xcolor}
\usepackage{titlesec}
\usepackage{enumitem}
\usepackage{array}
\usepackage{booktabs}
\usepackage{tabularx}
\usepackage{parskip}
\usepackage{fancyhdr}
\usepackage{lastpage}
\usepackage{needspace}
\usepackage{setspace}
\usepackage[
  colorlinks=true,
  linkcolor=navyblue,
  urlcolor=navyblue,
  pdfauthor={Muhebullah Karimzada},
  pdftitle={Teaching Dossier — Muhebullah Karimzada}
]{hyperref}

% ----- Colours -----
\definecolor{navyblue}{RGB}{0,32,96}
\definecolor{accentblue}{RGB}{30,80,160}
\definecolor{lightgray}{RGB}{130,130,130}
\definecolor{rulecolor}{RGB}{0,32,96}
\definecolor{boxbg}{RGB}{245,247,252}

% ----- Section formatting -----
\titleformat{\section}
  {\large\bfseries\color{navyblue}}
  {}{0em}{}
  [\vspace{1pt}{\color{navyblue}\titlerule[0.6pt]}]
\titlespacing{\section}{0pt}{18pt}{8pt}

\titleformat{\subsection}
  {\normalsize\bfseries\color{accentblue}}
  {}{0em}{}
\titlespacing{\subsection}{0pt}{10pt}{3pt}

% ----- Page style -----
\pagestyle{fancy}
\fancyhf{}
\renewcommand{\headrulewidth}{0pt}
\fancyhead[L]{\small\color{lightgray}\textit{Muhebullah Karimzada --- Teaching Dossier}}
\fancyhead[R]{\small\color{lightgray}\textit{September 2026}}
\fancyfoot[C]{\small\color{lightgray}\thepage\ of \pageref{LastPage}}

% ----- List settings -----
\setlist[itemize]{
  leftmargin=1.5em,
  itemsep=2pt,
  topsep=3pt,
  parsep=0pt
}
\setlist[enumerate]{
  leftmargin=1.5em,
  itemsep=2pt,
  topsep=3pt,
  parsep=0pt
}

% ----- Paragraph spacing -----
\setlength{\parskip}{5pt}
\setlength{\parindent}{0pt}

% ============================================================
\begin{document}
\thispagestyle{empty}

% ----- HEADER -----
\begin{center}
  {\LARGE\bfseries\color{navyblue} Muhebullah Karimzada, Ph.D.}\\[4pt]
  {\normalsize Assistant Professor, School of Economics}\\[2pt]
  {\normalsize University of Northern British Columbia, Prince George, British Columbia}\\[6pt]
  \small
  \href{mailto:muhebullah.karimzada@unbc.ca}{muhebullah.karimzada@unbc.ca}
\end{center}

\vspace{2pt}
\noindent{\color{navyblue}\rule{\textwidth}{0.8pt}}

\begin{center}
  {\large\bfseries\color{navyblue} Teaching Dossier}
\end{center}

\noindent{\color{navyblue}\rule{\textwidth}{0.4pt}}
\vspace{8pt}

% ----- TEACHING PROFILE AT A GLANCE -----
\noindent
\fcolorbox{navyblue}{boxbg}{\parbox{\dimexpr\linewidth-2\fboxsep-2\fboxrule\relax}{%
\vspace{6pt}
\centering\small
\begin{tabular}{@{}c@{\hspace{22pt}}c@{\hspace{22pt}}c@{\hspace{22pt}}c@{}}
  \textbf{9 years} & \textbf{4 universities} & \textbf{14 course preparations} & \textbf{2 original textbooks} \\
  university teaching & across Canada & principles to graduate level & in development \\
\end{tabular}
\vspace{6pt}
}}

\vspace{10pt}

% ============================================================
% TABLE OF CONTENTS
% ============================================================
\noindent\small
\begin{tabular}{@{}p{0.09\linewidth}@{\hspace{4pt}}p{0.86\linewidth}@{}}
  \textbf{I.}   & Teaching Philosophy and Approach \\[1pt]
  \textbf{II.}  & Teaching Record \\[1pt]
  \textbf{III.} & Curriculum Development and Course Design \\[1pt]
  \textbf{IV.}  & Evidence of Teaching Effectiveness \\[1pt]
  \textbf{V.}   & Mentorship and Student Engagement \\[1pt]
  \textbf{VI.}  & Professional Development and Service \\[1pt]
  \textbf{VII.} & Future Teaching Directions \\[1pt]
              & Appendix: Supporting Documentation \\
\end{tabular}
\normalsize

\vspace{12pt}

% ============================================================
\section{Part I: Teaching Philosophy and Approach}
% ============================================================

\subsection{Core Philosophy}

Teaching economics is, to me, an act of storytelling through theory and data. I approach each course as an opportunity to help students see how economic models interpret the world around them. My goal is not only to cultivate an understanding of economic principles but to develop an economist's mindset: analytical, empirical, and genuinely curious about why things are the way they are.

Economics is often taught as a collection of models and diagrams. My approach is different: I treat models as lenses that clarify real phenomena. A student who can look at a housing affordability crisis, a central bank rate decision, or a trade dispute and ask \textit{what forces are at work here?} has internalized economics in a way that endures far beyond the examination. I design my courses to build that instinct.

At a practical level, my teaching rests on three pillars: \textbf{conceptual clarity}, \textbf{application to Canadian evidence}, and \textbf{active engagement}. Running through all three is a commitment to treating students as capable analysts who are developing genuine understanding --- not as recipients of information to be memorized and reproduced.

\subsection{Conceptual Clarity}

I introduce intuition before formalism. When teaching aggregate demand and supply in ECON~101, I ask students first to reason through how a sudden rise in oil prices might affect output and employment in Prince George. By the time we reach the formal AD-AS diagram, students already understand the forces at work; the model confirms what they have already reasoned through. This sequence --- intuition, then formalism, then verification --- helps students trust their own economic reasoning rather than simply memorizing a graph.

In upper-division courses, this principle extends to mathematical structure. In ECON~320 (Introduction to Mathematical Economics), I introduce the concept of a derivative not as a calculus operation but as the answer to the question: \textit{how does a small change in this variable affect the outcome?} Students encounter this idea concretely before seeing the formal definition. By the time we reach constrained optimization and the envelope theorem, the mathematics feels like a language for expressing ideas they already hold, rather than a foreign formalism imposed from outside.

\subsection{Application and Canadian Evidence}

Economics must be \textit{lived}, not memorized. I integrate Canadian and international data into every course, drawing from Statistics Canada, the Bank of Canada, the IMF, and the World Bank. Students use real datasets to visualize business cycles, inflation, and labour market dynamics. Assignments include short empirical projects, policy briefs, and data-interpretation exercises designed to develop quantitative literacy alongside analytical writing.

Specific applications across my courses include:

\begin{itemize}
  \item \textbf{ECON 101 (Principles of Macroeconomics):} Students track the Bank of Canada's interest rate decisions alongside CPI data across the semester. Rate announcements become live case studies in aggregate demand management. In both Fall 2025 and Fall 2026, rate decisions landed during the term and were incorporated into lecture that week.

  \item \textbf{ECON 308 (International Economic Relations):} Lectures on the gravity model of trade use Canada--US and Canada--Japan trade flows; the effects of CPTPP and CUSMA are analyzed with real export data. Dedicated ``Canadian applications'' sections in every lecture connect trade theory to Dutch Disease, commodity terms-of-trade cycles, and Canada's exposure to U.S.\ trade policy.

  \item \textbf{ECON 300 (Labour Economics):} Students use Statistics Canada data to examine wage inequality trends, employment by industry, and the economics of immigration. Assignments require students to read and interpret published empirical studies on Canadian labour markets.

  \item \textbf{ECON 204 (Contemporary Economic Issues):} Each chapter is anchored in a Canadian policy debate --- the Bank of Canada's 2022--24 tightening cycle, housing affordability in Prince George and Vancouver, income inequality including Indigenous economic inequality in Canada, and the labour-market effects of AI adoption. A Canadian-authored textbook is currently under development for this course because no commercially available text addresses these topics from a Canadian institutional perspective.

  \item \textbf{ECON 311 (Behavioural Economics):} Canadian applications include automatic RRSP enrollment design, default rules in organ donation frameworks across provinces, price anchoring in regional housing markets, and behavioural dimensions of public health communication.

  \item \textbf{ECON 320 (Mathematical Economics):} Optimization problems use production functions and utility functions calibrated to realistic economic parameters. Students see that the mathematical tools they are learning describe real economic decisions.
\end{itemize}

\subsection{Active Engagement}

I foster a classroom environment where students reason through problems rather than receive conclusions. Specific engagement strategies I use regularly include:

\begin{itemize}
  \item \textbf{Structured problem-solving:} In-class quizzes are held at regular intervals --- typically six per semester, with the lowest score dropped. This structure rewards consistent engagement and reduces the high-stakes pressure of a single midterm.

  \item \textbf{Policy simulation:} In ECON 101, student groups periodically take the role of Bank of Canada analysts. Given current CPI, unemployment, and GDP data, they must justify a rate decision to the class. This exercise develops both quantitative reasoning and oral communication skills.

  \item \textbf{Discussion-based framing:} I open lectures with a question drawn from a recent news event or data release. In ECON 308, I have used tariff announcements, WTO dispute rulings, and commodity price shocks as entry points into theoretical frameworks. Students who engage with the question first are better prepared to appreciate the model.

  \item \textbf{Classroom experimentation:} In ECON 311, several lectures are structured as real-time experiments in which students make choices under conditions designed to elicit cognitive biases. Students analyze their own responses and compare them with the predictions of both standard and behavioural models --- making the theoretical content experientially credible rather than merely asserted.

  \item \textbf{Custom digital simulations:} I have developed Excel and Python-based simulations for topics including the money multiplier, the IS-LM model, the output gap, and exchange rate dynamics. These tools allow students to manipulate parameters and observe model outcomes directly, bridging the gap between theory and intuition.
\end{itemize}

\subsection{Assessment Philosophy}

My assessment approach is grounded in the principle that learning is incremental and that assessment should scaffold rather than merely measure that process. Several design choices follow from this commitment.

I distribute assessment across the term. Rather than concentrating weight in a single high-stakes final, I use a structure of regular in-class quizzes --- typically six, with the lowest dropped --- alongside two midterms and a final examination. This structure rewards consistent engagement and allows students to recover from a difficult week without catastrophic consequences for their grade.

I provide substantive written feedback on all assessed work. Students receive not just a grade but a brief explanation of where their reasoning was sound, where it broke down, and what they should do differently on the next assessment. Students who receive this targeted feedback are, in my experience, significantly more likely to engage with the corrected material before the next quiz or midterm.

I align assessment tasks tightly with instructional emphasis. If a lecture sequence has built toward connecting a model to Canadian data, the associated assessment asks students to perform exactly that connection on a new example. If a problem-solving session has walked through a constrained optimization problem step by step, the quiz requires students to solve a parallel problem independently. Assessment and instruction should tell the same story.

I maintain high and consistent standards across all student groups. Inclusive course design does not mean reduced expectations; it means ensuring that every student has a genuine opportunity to meet them. Supplementary materials, worked examples, and structured practice sets are provided precisely so that no student faces an assessment having lacked the opportunity to prepare adequately.

\subsection{Digital Tools and Pedagogical Technology}

Technology is most valuable in the classroom when it makes invisible processes visible. I have developed several digital tools for this purpose.

\textbf{Excel-based simulations.} For ECON 101, I built an interactive Bank of Canada policy simulation in which students observe how changes in the overnight rate propagate through the IS-LM model, shifting aggregate demand and affecting output and inflation. For ECON 308, I developed an Excel-based gravity model of trade in which students input bilateral trade data for Canada and its partners --- calibrated with actual World Bank data --- and compare model predictions with observed trade flows.

\textbf{Python simulations.} For intermediate and upper-level courses, I have developed Python notebooks covering exchange rate dynamics and the output gap. These notebooks allow students to change model parameters and observe simulated paths of key macroeconomic variables, connecting the formal model structure to what a time series of data would look like.

\textbf{R labs for econometrics.} Each chapter of the ECON 621 course notes is accompanied by a complete R lab with fully reproducible code. Students run OLS, IV, and difference-in-differences regressions on Canadian and international datasets. The labs are designed so that a student with no prior programming experience can execute them from the first session.

\textbf{Moodle and pre-lecture access.} All lecture notes, worked solutions, and practice problem sets are posted on Moodle before each lecture. Pre-posting materials allows students to engage with course content at their own pace and substantially improves the quality of in-class discussion, because students arrive having already encountered the vocabulary and key questions of the day's session.

\subsection{Inclusivity, Accessibility, and the UNBC Context}

UNBC serves a student population with distinctive needs and strengths. Many students come from Northern BC communities --- some from rural or remote areas far from major urban universities. A significant proportion are Indigenous students who bring different educational experiences and cultural frameworks into the economics classroom. A meaningful share are the first in their families to pursue a university degree.

I design my courses with this context explicitly in mind. In ECON 204, I include a dedicated section on Indigenous economic inequality in Canada: the Gini coefficient across Indigenous and non-Indigenous populations, the economics of reserve land systems, and the structural underfunding of Indigenous community infrastructure. This content is not appended as a ``special topic''; it is integrated into the analytical framework of the chapter and addressed with the same theoretical tools used throughout the course.

In all courses, worked examples and policy illustrations incorporate Northern BC contexts: Prince George's housing market, the regional labour market effects of the natural resource sector, and the economic implications of Northern infrastructure. Students from these communities see their own economic environment reflected in the material they are learning.

Office hours are structured to be genuinely accessible. I hold multiple slots each week and meet individually by appointment. I am aware that first-generation university students may be less inclined to seek out a professor proactively, and I actively invite students during lecture and through course communications. All lecture notes, worked solutions, and practice problem sets are posted before each class so that students can engage on their own schedule.

I believe that meeting students where they are, while maintaining rigorous academic standards, is not a contradiction. The two reinforce each other when course design is thoughtful.

% ============================================================
\section{Part II: Teaching Record}
% ============================================================

I have held teaching appointments at four Canadian universities over nine years. All courses listed were taught as the instructor of record with full responsibility for course design, lectures, assessments, and grading.

\subsection{University of Northern British Columbia --- Assistant Professor (January 2025 -- Present)}

\vspace{4pt}

\begin{tabular}{@{}>{\raggedright\arraybackslash}p{0.12\linewidth}@{\hspace{4pt}}>{\raggedright\arraybackslash}p{0.12\linewidth}@{\hspace{4pt}}>{\raggedright\arraybackslash}p{0.40\linewidth}@{\hspace{4pt}}>{\raggedright\arraybackslash}p{0.26\linewidth}@{}}
\toprule
\textbf{Term} & \textbf{Course} & \textbf{Title} & \textbf{Note} \\
\midrule
Winter 2025 & ECON 101 & Principles of Macroeconomics     & 34 students; new prep. \\
Winter 2025 & ECON 308 & International Economic Relations  & 6 students; new prep. \\
Winter 2025 & ECON 320 & Intro.\ to Mathematical Economics & 4 students; new prep. \\
Fall 2025   & ECON 101 & Principles of Macroeconomics     & 63 students \\
Fall 2025   & ECON 204 & Contemporary Economic Issues      & 26 students; new prep. \\
Fall 2025   & ECON 220 & Global Economic Shifts            & 19 students; new prep. \\
Winter 2026 & ECON 101 & Principles of Macroeconomics     & In person \\
Winter 2026 & ECON 204 & Contemporary Economic Issues      & In person \\
Winter 2026 & ECON 300 & Labour Economics                  & In person; new prep. \\
Fall 2026   & ECON 101 & Principles of Macroeconomics     & In person; current \\
Fall 2026   & ECON 308 & International Economic Relations  & In person; major redesign \\
Fall 2026   & ECON 311 & Behavioural Economics             & In person; new prep. \\
Fall 2026   & ECON 320 & Intro.\ to Mathematical Economics & In person \\
\bottomrule
\end{tabular}

\vspace{6pt}

\textit{Note:} ``New prep.'' indicates a course designed and delivered for the first time, including original lecture materials, assessments, and a course outline. Enrollment figures for recent terms are available upon request.

\subsection{University of Northern British Columbia --- Sessional Lecturer (Fall 2022, Fall 2023)}

Independently taught introductory and upper-division economics courses --- including Principles of Macroeconomics, Contemporary Economic Issues, and Global Economic Shifts --- prior to the full-time appointment. These offerings established my familiarity with UNBC's student population and curriculum context before joining the faculty full-time.

\subsection{University of Waterloo --- Sessional Lecturer (January 2018 -- May 2025)}

\begin{tabular}{@{}p{0.55\linewidth}@{\hspace{6pt}}p{0.36\linewidth}@{}}
  Introduction to Macroeconomics & Second-year undergraduate \\
  Money and Banking I            & Second-year undergraduate \\
  International Finance          & Upper-division undergraduate \\
  Macroeconomics Theory III      & Upper-division undergraduate \\
\end{tabular}

\vspace{4pt}
Courses at Waterloo required coverage at a higher mathematical level than standard introductory offerings, reflecting the rigorous macroeconomic and monetary theory tradition of that department. My preparation and delivery of these courses directly informs the upper-division material I teach at UNBC.

\subsection{McMaster University --- Sessional Lecturer (2017, 2018, 2022, 2023)}

\begin{tabular}{@{}p{0.55\linewidth}@{\hspace{6pt}}p{0.36\linewidth}@{}}
  Introduction to Macroeconomics     & First-year undergraduate \\
  International Monetary Economics   & Upper-division undergraduate \\
  Intermediate Macroeconomics II     & Third-year undergraduate \\
\end{tabular}

\vspace{4pt}
The McMaster introductory course enrolled several hundred students per section. Managing large-section teaching --- maintaining engagement, designing scalable assessment, and providing efficient feedback --- is a skill I developed there and continue to draw on.

\subsection{University of Toronto --- Sessional Lecturer (2023--2024)}

\begin{tabular}{@{}p{0.55\linewidth}@{\hspace{6pt}}p{0.36\linewidth}@{}}
  Topics in Macroeconomic Theory & Upper-division undergraduate \\
\end{tabular}

\vspace{4pt}
Classroom teaching at the University of Toronto was formally evaluated by a peer observer during this appointment. See Part IV for the full evaluation.

\subsection{Teaching Coverage by Area}

\vspace{4pt}

\begin{tabular}{@{}>{\raggedright\arraybackslash}p{0.29\linewidth}@{\hspace{6pt}}>{\raggedright\arraybackslash}p{0.63\linewidth}@{}}
\textbf{Macroeconomics}        & Principles; Intermediate Macroeconomics II; Macroeconomics Theory III; Topics in Macroeconomic Theory \\[3pt]
\textbf{Monetary Economics}    & Money and Banking I; International Monetary Economics; International Finance \\[3pt]
\textbf{International Econ.}   & International Economic Relations (Trade); International Finance; International Monetary Economics; Global Economic Shifts \\[3pt]
\textbf{Labour Economics}      & Labour Economics (ECON 300) \\[3pt]
\textbf{Mathematical Econ.}    & Introduction to Mathematical Economics (ECON 320) \\[3pt]
\textbf{Behavioural Econ.}     & Behavioural Economics (ECON 311) \\[3pt]
\textbf{Contemporary Issues}   & Contemporary Economic Issues (ECON 204) \\
\end{tabular}

% ============================================================
\section{Part III: Curriculum Development and Course Design}
% ============================================================

\subsection{Overview}

Since joining UNBC as a full-time faculty member in January 2025, I have developed nine new course preparations and substantially redesigned two existing ones --- all within twenty months. This pace reflects both the genuine curriculum needs of a growing economics program and my belief that original, carefully designed instructional materials are a precondition for excellent teaching. I do not adapt materials from other institutions when I can develop materials specifically designed for the students I teach.

\subsection{New Course Packages Developed at UNBC}

\textbf{Online Course Packages (2024--25).} In preparation for my first year of full-time teaching, I designed three complete online course packages --- ECON 101, ECON 204, and ECON 220 --- for asynchronous delivery. Each package included a full set of lecture slide decks, narrated digital lectures, structured practice problem sets with full solutions, tutorial exercises, six quizzes, two midterm examinations, and a final examination. The courses were developed from scratch; no prior UNBC materials were adapted. Designing for asynchronous delivery imposes a useful discipline: problems must be perfectly stated, solutions fully worked, and explanations capable of standing on their own without the supplementary clarifications that an in-person lecture naturally provides.

\textbf{ECON 308 (International Economic Relations) --- Initial Design, Winter 2025.} The first offering of ECON 308 under my instruction used Krugman, Obstfeld, and Melitz as the primary text. I developed a complete set of lecture notes, four problem sets with solutions, two midterms, and a final examination. This initial offering informed the comprehensive redesign for Fall 2026 described below.

\textbf{ECON 320 (Introduction to Mathematical Economics) --- Initial Design, Winter 2025.} This course covers comparative statics, differentiation rules, optimization, exponential and logarithmic functions, constrained optimization, and an introduction to dynamics and differential equations. I developed the full course structure around Chiang and Wainwright's \textit{Fundamental Methods of Mathematical Economics} (4th edition), with six in-class quizzes and two midterms. A distinguishing feature of the course is its emphasis on economic interpretation at every step: every mathematical result is immediately connected to what it means for an economic agent's decision.

\textbf{ECON 300 (Labour Economics) --- New Design, Winter 2026.} I developed this course using Benjamin, Gunderson, Lemieux, Schirle, and Riddell's \textit{Labour Market Economics} (9th edition), with a coverage sequence emphasizing labour supply and demand, wages, human capital, unemployment, and immigration. The textbook was supplemented with Statistics Canada data exercises and readings from contemporary Canadian empirical research on Northern and Indigenous labour markets. The assessment design includes five quizzes, one empirical assignment, two midterms, and a final examination.

\textbf{ECON 311 (Behavioural Economics) --- New Preparation, Fall 2026.} This course examines how systematic psychological mechanisms --- loss aversion, present bias, overconfidence, social preferences, and framing effects --- produce predictable departures from standard rational-choice theory, and how those departures shape individual decisions and policy outcomes. The course is organized around three themes: the psychological mechanisms themselves and their theoretical foundations; the laboratory and field experimental evidence that documents them; and the policy implications, including the design of choice architectures and nudges.

The primary text is Angner's \textit{A Course in Behavioral Economics}, supplemented by readings from Kahneman, Thaler, and Sunstein, and by key field experiment papers. Theoretical topics include prospect theory and loss aversion, present bias and hyperbolic discounting, social preferences and fairness norms, overconfidence and the representativeness heuristic. Each topic is paired with an empirical application: retirement savings and automatic enrollment (present bias); organ donation default rules across Canadian provinces (nudge design); price anchoring in housing markets (framing effects); and vaccination take-up and public health messaging (social preferences).

A distinctive pedagogical feature of ECON 311 is a sequence of in-class experiments in which students make choices under conditions designed to elicit predictable cognitive biases. Students observe their own decisions, compare them with the predictions of both standard and behavioural models, and analyze the divergence. This approach --- developed in the tradition of Kahneman and Tversky's original framing experiments --- gives students firsthand insight into the mechanisms they are studying and makes the theoretical content experientially credible. Assessment includes six short analytical assignments, a midterm, and a final examination. Each assignment asks students to identify a behavioural economic phenomenon in a recent Canadian news event, describe the underlying cognitive mechanism, and propose a policy response --- developing analytical writing alongside the habit of reading economic news critically.

\subsection{Major Course Redesigns}

\textbf{ECON 101 --- In-Person Redesign, Fall 2025.} The transition from online to in-person delivery required a complete pedagogical restructuring. I introduced active-learning components including in-class problem sessions, small-group discussion exercises, and real-time data analysis segments using live Bank of Canada data. The assessment structure was revised to include six in-class quizzes alongside two midterms and a final examination, increasing continuous assessment and reducing the weight of the final. ECON 101 enrollment grew from 34 students in Winter 2025 (online) to 63 in Fall 2025 (in-person), and the in-person format has continued into subsequent terms.

\textbf{ECON 308 --- Comprehensive Redesign, Fall 2026.} After reviewing student performance and course feedback from Winter 2025, I undertook a comprehensive redesign of ECON 308. The redesign involved:

\begin{itemize}
  \item Adopting the 12th edition of Krugman, Obstfeld, and Melitz, with a chapter sequencing that introduces trade theory before trade policy.
  \item Creating original lecture slides in \LaTeX/Beamer with custom TikZ-typeset diagrams for all core models: Ricardian, Specific Factors, Heckscher-Ohlin, Standard Trade Model, and External Economies.
  \item Developing a pedagogical sequence that integrates trade theory with open-economy macroeconomics at each step, connecting comparative advantage to exchange rates, the current account, and Canadian macroeconomic outcomes.
  \item Adding dedicated ``Canadian applications'' sections in every lecture, covering Dutch Disease, commodity terms-of-trade cycles, CUSMA wage and rules-of-origin provisions, and Canada's exposure to U.S.\ trade policy.
  \item Building new worked examples with numerical solutions drawn from Canadian trade data.
\end{itemize}

\subsection{Instructional Materials in Progress}

\textbf{\textit{Contemporary Economic Issues: A Canadian Perspective.}} I am developing an original seven-chapter \LaTeX{} textbook for ECON 204 / INTS 211. No commercially available text for this course type addresses Canadian institutions, data, or policy debates in depth; the market is dominated by U.S.-centric introductory texts that leave Canadian students reading examples with no connection to their own economic context. The book covers: (1)~introduction to economic reasoning and Canadian economic institutions; (2)~inflation, cost of living, and monetary policy; (3)~income inequality, including a dedicated chapter section on Indigenous economic inequality in Canada; (4)~housing affordability; (5)~labour markets and AI; (6)~technology and productivity; and (7)~climate change and environmental economics. Each chapter is anchored to a live Canadian policy debate, uses Statistics Canada or Bank of Canada data, and includes discussion questions and policy exercises. The book is designed for adoption in Winter 2027.

\textbf{\textit{Applied Econometrics} --- Course Notes for ECON 621.} I have developed thirteen chapters of original course notes for ECON 621 (Applied Econometrics), a new cross-listed senior undergraduate and graduate course planned for Winter 2027. The notes cover potential outcomes and causal inference; the classical linear regression model; OLS algebra; finite-sample and asymptotic properties; the Frisch-Waugh-Lovell theorem; instrumental variables and two-stage least squares; LATE; panel data and fixed effects; difference-in-differences; matching and propensity score methods; and regression discontinuity design. Each chapter includes worked economic examples, analytical exercises, and a complete R lab with reproducible code. The notes are designed to be accessible to senior undergraduates without prior exposure to matrix algebra, while retaining the formal rigour appropriate for graduate-level instruction.

% ============================================================
\section{Part IV: Evidence of Teaching Effectiveness}
% ============================================================

\subsection{Student Course Evaluations}

Teaching evaluations have been administered across my appointments at McMaster University, the University of Waterloo, and the University of Toronto. Official copies are included in the appendix.

\textbf{McMaster University --- Full Survey Instructor Report (ECON 1BB3: Introduction to Macroeconomics).} A standardized instructor survey was administered for this course, covering overall quality of instruction, course organization, availability for student questions, and clarity of explanation. Results are available in the appendix.

\textbf{University of Waterloo --- Official Student Course Perception Results (ECON 102: Money and Banking I).} Waterloo's official course perception instrument was administered for this course, producing a standardized report. Results are included in the appendix.

\textbf{University of Toronto --- Student Course Evaluations.} Student evaluations from Topics in Macroeconomic Theory are included in the appendix.

Representative student observations from across these evaluations include:

\begin{itemize}
  \item ``This is the first economics course where I felt like I genuinely understood \textit{why} the model works, not just how to draw the diagram.''
  \item ``The instructor connects the theory to things happening in Canada right now. It makes the course feel relevant.''
  \item ``I appreciated the quizzes --- they kept me engaged week to week and I felt much more prepared for the midterm.''
  \item ``The instructor is very accessible and responds to questions thoughtfully. He takes the time to make sure you understand before moving on.''
  \item ``I've never had an economics professor who explains the intuition this clearly. Most just derive the equation; he makes sure you understand what the equation is actually saying.''
  \item ``The worked examples in lecture were excellent. I could follow them step by step and then immediately try the practice problems.''
  \item ``The course made me realize I actually want to study economics further. I came in thinking it was just graphs and I'm leaving having genuinely enjoyed it.''
\end{itemize}

At UNBC, student course perception surveys are conducted through the standard institutional process. Results from Winter 2025, Fall 2025, and subsequent terms are available upon request.

\subsection{Responsiveness to Student Feedback}

I treat student evaluations as data. After each term, I review feedback systematically, identify patterns, and make specific, documented changes to course design and delivery. This reflective practice has directly informed several of the revisions described in Part III.

Following Winter 2025 feedback on ECON 308, students noted that coverage pace was too fast in the second half of the term and that worked numerical examples were sparse in the trade-policy sections. Both observations were incorporated into the Fall 2026 redesign: chapter sequencing was adjusted to give more time to trade policy, and every policy lecture now includes at least one fully worked numerical example drawn from Canadian trade data.

Following the first in-person offering of ECON 101 in Fall 2025, students indicated strong satisfaction with the quiz structure but requested more explicit connections between lecture content and examination expectations. I responded by introducing a brief ``exam connection'' segment at the end of each lecture, identifying which analytical skills from the day's material would appear on assessments and in what form. Student preparation for midterms improved noticeably in subsequent terms.

I share selected evaluation feedback with students at the beginning of the following term, explaining what I changed in response to their input. This transparency reinforces that student feedback is taken seriously and builds a culture of mutual accountability in the classroom.

\subsection{Peer Teaching Evaluation}

During my appointment at the University of Toronto, my teaching was formally observed and evaluated by \textbf{Dr.\ Laurent Cavenaile} (Department of Economics, University of Toronto). The evaluation included both a classroom observation and a review of course documents (syllabus, assessments, and lecture materials). A copy of the written peer evaluation report is included in the appendix.

Dr.\ Cavenaile's assessment addressed the clarity and organization of lecture delivery, the quality of course materials, the appropriateness of assessments, and the overall pedagogical approach. His observation noted particular strengths in the connection between theoretical frameworks and applied examples, and in the management of student engagement during lecture.

\subsection{Evidence of Student Learning and Outcomes}

Student performance data across my courses indicate strong and consistent engagement with course material. At UNBC, the continuous assessment structure --- six quizzes, two midterms, final examination --- produces a clear pattern: students' performance from the first midterm to the second typically improves, suggesting that the quiz-driven feedback loop is functioning as designed. Students who are engaged week to week are better prepared when the high-stakes assessments arrive.

ECON 101 enrollment grew from 34 students in Winter 2025 to 63 in Fall 2025, reflecting both program growth and word-of-mouth engagement with the course structure. I regard enrollment growth not as a primary metric of teaching effectiveness but as indirect evidence that students find genuine value in the course.

Beyond examination performance, a number of students I have mentored have been admitted to highly competitive graduate programs in economics in Canada and internationally. One former student contacted me after receiving multiple offers of admission and expressed that my teaching had been formative in developing their interest in and preparation for graduate-level work. A copy of that correspondence is included in the appendix.

% ============================================================
\section{Part V: Mentorship and Student Engagement}
% ============================================================

\subsection{Academic Advising and Graduate School Preparation}

I maintain an open-door advising policy and hold multiple office hours each week. Students visit to discuss course material but also to talk through graduate school preparation, career directions, and research interests. I encourage students who are curious about research to engage directly with working papers and empirical literature, not only textbooks.

I have supported students in preparing graduate school applications, sharpening research statements, and identifying programs appropriate to their preparation and goals. I am attentive to the particular challenges faced by students at UNBC: many are the first in their families to consider graduate education and may underestimate the competitiveness of their own preparation. I take it as part of my responsibility to help students see themselves as capable of advanced work when the evidence supports that view.

I have written letters of recommendation for students pursuing graduate programs in economics, public policy, and business. I take this responsibility seriously: a strong letter is specific, analytical, and credible. I write letters that will be read as genuine professional assessments, not boilerplate endorsements, and I engage with each student's record carefully before agreeing to write.

For students seriously considering graduate economics, I provide supplementary guidance on the mathematical and econometric preparation that graduate programs expect --- preparation that is not always visible from the undergraduate curriculum, and where early guidance can substantially affect outcomes.

\subsection{Student Research Projects and Research-Adjacent Experiences}

Beyond formal coursework, I create research-adjacent experiences through both course design and individual mentorship. In ECON 300, a substantial assignment requires students to identify a research question in Canadian labour economics, locate the relevant empirical literature, describe the empirical approach they would use, and interpret the published evidence. This assignment approximates the structure of a policy research project and serves as an early introduction to the research process for students who have not yet encountered formal research methods.

I have directed students toward working papers and empirical studies relevant to their interests, helped identify Statistics Canada data sources for independent projects, and provided informal mentorship on project design. At UNBC, where the formal infrastructure for undergraduate research in economics is still developing, I try to create these opportunities through the course and advising relationship.

I use my own research experience --- in macroeconomics, DSGE modelling, and international economics --- to show students what economists do beyond the classroom. I discuss the process of formulating a research question, building a formal model, and confronting the data. I want students to understand economics not only as a body of knowledge but as an ongoing method of inquiry, and to see the models in their textbooks as live investigations rather than settled facts.

\subsection{The UNBC Student Population: Differentiated Support and Northern Context}

UNBC serves a diverse student population that includes Indigenous students, students from rural and remote Northern BC communities, and many who are the first in their families to attend university. This is not a demographic footnote; it is the central context in which I teach.

Quantitative courses such as ECON~320 include supplementary material on foundational mathematics for students who arrive with gaps in that preparation --- without signalling a lower standard. I present these materials as resources for building strength, not as remediation. In ECON~204 and ECON~101, additional scaffolding on quantitative reasoning is embedded in the course structure rather than offered as an afterthought.

I provide structured feedback on every assessment --- not only a grade but specific written notes indicating where reasoning was correct, where it broke down, and what the student should do differently. Students who receive targeted feedback engage more seriously with the corrected version of their thinking. I regard detailed feedback as one of the highest-value inputs I can provide, and I allocate time accordingly.

I am also attentive to how the lived experience of Northern and Indigenous economic conditions intersects with the economics curriculum. When students from communities affected by resource extraction cycles, land rights questions, or Indigenous infrastructure gaps encounter course content that names and analyzes those conditions analytically, the economics becomes real in a way that no hypothetical example can replicate. This is a distinctive strength of teaching at UNBC, and I design my courses to use it.

% ============================================================
\section{Part VI: Professional Development and Service}
% ============================================================

\subsection{Teaching-Related Professional Development}

\begin{tabular}{@{}p{0.62\linewidth}@{\hspace{6pt}}p{0.26\linewidth}@{}}
  ``How Can We Effectively Use AI for Research,'' Chalmers University of Technology (online) & September 2025 \\[2pt]
  Research Data Centre (RDC) Orientation, Statistics Canada / UNBC                           & October 2025  \\[2pt]
  Faculty and Staff Orientation, UNBC                                                         & October 2025  \\[2pt]
  ECON Development Forum, School of Economics, UNBC                                          & November 2025 \\[2pt]
  School of Economics Seminar Series, UNBC                                                   & December 2025 \\
\end{tabular}

\vspace{6pt}

Access to Statistics Canada's Research Data Centre (RDC) at UNBC directly supports the empirical dimension of my teaching. The RDC orientation in October 2025 provided formal training in the data governance and access procedures that govern Statistics Canada's confidential microdata --- data I draw on in ECON 300 and plan to incorporate in ECON 621.

\subsection{Connecting Teaching to Research}

My active research program in macroeconomics and international economics directly informs my teaching. Familiarity with the open-economy macroeconomics literature allows me to draw on recent empirical work --- on exchange rate pass-through, trade elasticities, and global value chains --- that goes well beyond the textbook. My work in DSGE modelling gives me fluency with the Bank of Canada's policy framework that enriches the in-class treatment of rate decisions in ECON 101 and ECON 308 beyond what any textbook account can provide.

I use my research experience explicitly with students. I discuss the process of formulating a question, building a formal model, and confronting the data. I want students to understand that the textbook models they are studying are not settled truths but the best current answers to questions that researchers are still actively working on --- and that economics as a discipline is an invitation to participate in that inquiry.

\subsection{Academic Service Related to Teaching and Curriculum}

\textbf{Chair, Student Recruitment Committee, School of Economics, UNBC (2025--Present).}
I lead the committee responsible for promoting the economics program to prospective students. Responsibilities include coordinating outreach with local high schools, designing recruitment materials, and developing strategies to strengthen program enrolment. I have attended campus open-house events and prospective student information sessions, and have contributed to orientation programming for incoming economics students.

\textbf{Member, Curriculum Committee, School of Economics, UNBC (2025--Present).}
I participate in the review of course offerings, updating of program learning outcomes, and evaluation of new course proposals. This committee work connects directly to my curriculum development activities and ensures that the courses I develop are aligned with program-level objectives and UNBC's institutional learning outcomes.

% ============================================================
\section{Part VII: Future Teaching Directions}
% ============================================================

My teaching priorities over the next three years are to deepen the Canadian applied content of existing courses, complete the two original textbooks in progress, expand the upper-division and graduate curriculum, and contribute to the long-term development of the UNBC School of Economics.

\textbf{Courses I am prepared to develop or teach immediately:}

\begin{itemize}
  \item \textit{Applied Econometrics} (ECON 621) --- thirteen chapters of course notes are complete; first offering planned for Winter 2027. This course fills a significant gap in the UNBC curriculum by giving senior students and graduate students the empirical tools required for policy-relevant research.
  \item \textit{Money, Banking, and Financial Systems} --- drawing on my Waterloo teaching record in Money and Banking I and my active research in monetary economics.
  \item \textit{Intermediate Macroeconomics} --- I have taught this material at McMaster and am prepared to develop a Canadian-context version with updated data applications and a focus on the Bank of Canada's macroeconomic framework.
  \item \textit{Canadian Economic Policy} --- a policy-oriented course integrating macroeconomics, trade, and labour economics with a Canadian institutional focus, designed for upper-division students and non-specialist graduate students.
  \item \textit{Advanced Macroeconomic Theory} --- for students preparing for graduate study, drawing directly on my research in DSGE modelling and business cycle theory.
\end{itemize}

\textbf{Longer-term teaching and program goals:}

I aim to build ECON~204 into a flagship course for the economics program --- a course that draws non-majors into the discipline by making economics immediately relevant to their lives. The Canadian textbook under development is designed to support this goal. A well-designed introductory economics course at a regionally-focused university is, in a meaningful sense, an act of civic education: students who leave the course understanding how Canadian labour markets, monetary policy, and trade relations function are better-equipped citizens, regardless of their career direction.

I also intend to develop an applied econometrics track at the senior undergraduate level, giving UNBC students the quantitative tools they need to engage with empirical research and enter graduate programs or quantitatively-oriented careers. This track would combine ECON 621 with a sequence of upper-division courses in which students apply causal inference methods to Canadian data on topics of regional and national importance.

\textbf{Teaching and UNBC's institutional mission:}

UNBC occupies a distinctive position in Canadian higher education: a primarily undergraduate, regionally-focused university in Northern BC with a deep commitment to Northern and Indigenous communities. The economics program has an opportunity to develop a curriculum that speaks directly to this context --- courses that take the Northern BC economy seriously as an object of analysis, that apply economic theory to the questions facing Northern and Indigenous communities, and that prepare students for the careers and graduate programs available from this region.

My trajectory at UNBC --- developing original Canadian-context textbooks, redesigning courses around Canadian evidence, building the applied econometrics curriculum, and chairing student recruitment --- is a deliberate contribution to this program vision. My goal is to help build a program that is recognizably UNBC's: rigorous, Canadian, and genuinely connected to the communities it serves.

% ============================================================
\needspace{5\baselineskip}
\section{Appendix: Supporting Documentation}
% ============================================================

The following documents are included with this dossier or available upon request:

\begin{enumerate}
  \item \textbf{Official Survey Instructor Report} --- ECON 1BB3: Introduction to Macroeconomics, McMaster University.
  \item \textbf{Official Student Course Perception Results} --- ECON 102: Money and Banking I, University of Waterloo.
  \item \textbf{Student Course Evaluations} --- Topics in Macroeconomic Theory, University of Toronto.
  \item \textbf{Peer Teaching Evaluation} --- Classroom observation and course document review by Dr.\ Laurent Cavenaile, Department of Economics, University of Toronto.
  \item \textbf{Student correspondence} --- Email from former student regarding graduate school admission, forwarded with permission.
  \item \textbf{Sample course syllabi} --- ECON 320 (Introduction to Mathematical Economics, Fall 2026); ECON 300 (Labour Economics, Winter 2026); ECON 308 (International Economic Relations, Fall 2026); ECON 311 (Behavioural Economics, Fall 2026).
  \item \textbf{Sample assessment materials} --- Problem sets, quiz questions, and midterm examinations from ECON 101, ECON 308, and ECON 320 available upon request.
  \item \textbf{Sample instructional materials} --- Excerpt from \textit{Contemporary Economic Issues: A Canadian Perspective} (ECON 204 textbook in progress) and sample R lab from ECON 621 course notes available upon request.
  \item \textbf{UNBC student course perception results} --- Available upon request for Winter 2025, Fall 2025, and subsequent terms.
  \item \textbf{Curriculum Vitae} --- Full academic CV provided separately.
\end{enumerate}

\end{document}
